\documentclass[11pt]{article}

% ── Geometry ──────────────────────────────────────────────────────────
\usepackage[margin=1in]{geometry}

% ── Fonts & encoding ─────────────────────────────────────────────────
\usepackage[T1]{fontenc}
\usepackage{lmodern}
\usepackage{microtype}

% ── Math ──────────────────────────────────────────────────────────────
\usepackage{amsmath,amssymb,amsthm,mathtools}

% ── Graphics & tables ────────────────────────────────────────────────
\usepackage{graphicx}
\usepackage{booktabs}
\usepackage{multirow}
\usepackage{subcaption}
\usepackage{xcolor}

% ── Algorithms ────────────────────────────────────────────────────────
\usepackage[ruled,vlined]{algorithm2e}

% ── References ────────────────────────────────────────────────────────
\usepackage[colorlinks,citecolor=blue!60!black,linkcolor=blue!60!black,urlcolor=blue!60!black]{hyperref}
\usepackage[numbers,sort&compress]{natbib}

% ── Theorem environments ─────────────────────────────────────────────
\newtheorem{theorem}{Theorem}[section]
\newtheorem{lemma}[theorem]{Lemma}
\newtheorem{proposition}[theorem]{Proposition}
\newtheorem{corollary}[theorem]{Corollary}
\newtheorem{conjecture}[theorem]{Conjecture}
\theoremstyle{definition}
\newtheorem{definition}[theorem]{Definition}
\newtheorem{example}[theorem]{Example}
\theoremstyle{remark}
\newtheorem{remark}[theorem]{Remark}

% ── Notation ──────────────────────────────────────────────────────────
\newcommand{\R}{\mathbb{R}}
\newcommand{\Z}{\mathbb{Z}}
\newcommand{\N}{\mathbb{N}}
\newcommand{\cI}{\mathcal{I}}
\newcommand{\cB}{\mathcal{B}}
\newcommand{\cC}{\mathcal{C}}
\newcommand{\cM}{\mathcal{M}}
\newcommand{\cN}{\mathcal{N}}
\newcommand{\cS}{\mathcal{S}}
\newcommand{\Gr}{\mathrm{Gr}}
\newcommand{\relu}{\operatorname{ReLU}}
\DeclareMathOperator{\rank}{rank}
\DeclareMathOperator{\supp}{supp}

% ── Page-break control ──────────────────────────────────────────────
\widowpenalty=10000
\clubpenalty=10000
\displaywidowpenalty=10000
\usepackage{needspace}

% ── Document ──────────────────────────────────────────────────────────
\title{Positroid Structure in ReLU Networks:\\
  From Combinatorial Geometry to Matroid-Guided Pruning}

\author{
  Harrison Totty \\
  \texttt{harrison@totty.dev}
}

\date{}

\begin{document}
\maketitle

% ══════════════════════════════════════════════════════════════════════
\begin{abstract}
% ══════════════════════════════════════════════════════════════════════

We study the combinatorial structure of hyperplane arrangements
induced by single-hidden-layer ReLU networks with totally positive
(TP) weight matrices.  The activation pattern of such a network
defines an \emph{affine matroid} on the hidden neurons, and we
investigate when this matroid is a \emph{positroid}---a matroid
arising from the totally nonnegative Grassmannian.

We prove two theorems.  The \textbf{Contiguous-Implies-Positroid
(CIP) Theorem} states that if every non-basis of a rank-$k$ matroid
on $[n]$ is a cyclic interval, then the matroid is a positroid.  The
\textbf{Contiguous-Support Positroid Theorem (CSPT)} generalizes
this: if the union of all elements appearing in non-bases forms a
cyclic interval of rank less than~$k$, the matroid is a positroid.
These theorems explain the empirical observation that trained
networks produce exclusively positroid activation matroids (800+
trials with TP-constrained weights, zero exceptions; unconstrained
controls are discussed in Section~\ref{sec:conjecture}).

We show that the TP constraint is essential by constructing 12{,}642
counterexamples to the stronger conjecture that TP weights alone
suffice---deliberately chosen biases with ``spread'' non-basis
patterns break positroid structure, while ``contiguous'' patterns
never do.  Non-TP parameterizations (negated bidiagonal) produce
non-positroids at a 3\% rate, confirming that total positivity is
the operative mechanism.

Finally, the matroid's rank-deficiency partition identifies a
\textbf{safe pruning set}: removing 75\% of tail neurons causes zero
accuracy loss across all scales tested ($H = 50$ to $200$),
outperforming random pruning by 10--15 percentage points and strictly
outperforming magnitude, activation, and sensitivity baselines.

\end{abstract}


% ══════════════════════════════════════════════════════════════════════
\section{Introduction}\label{sec:intro}
% ══════════════════════════════════════════════════════════════════════

A single-hidden-layer ReLU network
$f(x) = W_2 \cdot \relu(W_1 x + b_1) + b_2$
with $H$ hidden neurons defines a hyperplane arrangement in input
space: the $i$-th neuron activates on the half-space
$\{x : w_i^\top x + b_i \geq 0\}$, where $w_i$ is the $i$-th row
of $W_1$.  The combinatorial structure of this arrangement---which
subsets of hyperplanes can simultaneously pass through a
point---is captured by an \emph{affine matroid} of rank at most $d+1$ on
the ground set $[H] = \{0, 1, \ldots, H{-}1\}$, where $d$ is the
input dimension.

\emph{Positroids} are a distinguished class of matroids arising from
the totally nonnegative Grassmannian $\Gr_+(k,n)$
\citep{postnikov2006,oh2011,knutson2013}.  They enjoy rich
combinatorial structure---Grassmann necklaces, decorated
permutations, plabic graphs---and appear throughout algebraic
combinatorics and mathematical physics
\citep{lukowski2023}.

In this paper, we establish a new connection between positroid
combinatorics and neural network geometry.  Our contributions are:

\begin{enumerate}
  \item \textbf{Empirical conjecture and disproof.}  We conjecture
    that TP weight matrices always produce positroid activation
    matroids.  Across 800+ training trials on multiple datasets, the
    conjecture holds without exception.  We then disprove the
    conjecture by constructing 12{,}642 counterexamples via
    deliberately adversarial bias choices
    (Section~\ref{sec:counterexamples}).

  \item \textbf{Two matroid theorems.}  The
    \emph{Contiguous-Implies-Positroid} (CIP) theorem
    (Theorem~\ref{thm:cip}) and the \emph{Contiguous-Support
    Positroid Theorem} (CSPT, Theorem~\ref{thm:cspt}) provide a
    complete explanation of the empirical dichotomy: networks whose
    non-bases are cyclic intervals are always positroids, while
    ``spread'' non-bases can break the structure
    (Section~\ref{sec:theorems}).

  \item \textbf{TP is essential.}  Non-TP parameterizations
    (negated bidiagonal) produce non-positroids at a $\sim$3\% rate,
    while TP parameterizations never do.  The causal chain is:
    $\text{TP} \to \text{contiguous support} \to
    \text{CSPT} \to \text{positroid}$
    (Section~\ref{sec:tp-essential}).

  \item \textbf{Matroid-guided pruning.}  The matroid's
    rank-deficiency partition identifies tail neurons that can be
    removed with zero accuracy loss.  At 75\% tail removal, this
    strategy achieves 0\% accuracy loss at every scale tested
    ($H = 50$ to $200$), outperforming random pruning by 10--15
    percentage points and strictly outperforming magnitude,
    activation, and sensitivity baselines
    (Section~\ref{sec:pruning}).
\end{enumerate}


\subsection{Related work}\label{sec:related}

\paragraph{Positroid theory.}
Positroids were introduced by Postnikov \citep{postnikov2006} as
matroids realizable in $\Gr_+(k,n)$, with characterizations via
Grassmann necklaces, decorated permutations, and plabic graphs.  Oh
\citep{oh2011} proved that positroids are exactly intersections of
cyclically shifted Schubert matroids.
Knutson, Lam, and Speyer \citep{knutson2013} studied positroid
varieties.  Łukowski, Parisi, and Williams \citep{lukowski2023}
connected positroid combinatorics to the amplituhedron program.

\paragraph{ReLU network geometry.}
The hyperplane arrangement perspective on ReLU networks has been
studied by several authors.  Zhang, Naitzat, and Lim
\citep{zhang2018} studied neural networks through tropical geometry.
Hanin and Rolnick \citep{hanin2019} counted linear regions.  The
connection between network architecture and arrangement combinatorics
remains largely unexplored from the matroid-theoretic viewpoint.

\paragraph{Implicit bias of gradient descent.}
Soudry et~al.\ \citep{soudry2018} showed that gradient descent on
separable data converges in direction to the max-margin classifier.
Gunasekar et~al.\ \citep{gunasekar2018} studied implicit bias toward
sparsity in the frequency domain for linear convolutional networks.  Our work reveals a new form of implicit bias:
\emph{combinatorial} implicit bias, where training induces a
specific matroid type (positroid) rather than a specific
norm or rank.

\paragraph{Neural network pruning.}
Magnitude pruning \citep{han2015} removes weights by absolute value.
Structured pruning methods identify entire neurons or channels for
removal.  Our approach is fundamentally different: we use the
matroid structure of the hyperplane arrangement to identify neurons
whose removal is \emph{provably} lossless with respect to the
arrangement's combinatorics.


% ══════════════════════════════════════════════════════════════════════
\section{Setup and the conjecture}\label{sec:setup}
% ══════════════════════════════════════════════════════════════════════

\subsection{Matroids and positroids}\label{sec:matroids}

\begin{definition}[Matroid]
  A \emph{matroid} $\cM = (E, \cB)$ is a finite ground set $E$ with
  a nonempty collection $\cB \subseteq 2^E$ of \emph{bases}---subsets
  of equal cardinality $k = \rank(\cM)$---satisfying the exchange
  axiom: for any $B_1, B_2 \in \cB$ and $x \in B_1 \setminus B_2$,
  there exists $y \in B_2 \setminus B_1$ such that
  $(B_1 \setminus \{x\}) \cup \{y\} \in \cB$.
\end{definition}

\begin{definition}[Linear and affine matroid]
  Given vectors $v_1, \ldots, v_m \in \R^n$, the \emph{linear
  matroid} has ground set $[m]$ and bases
  $\cB = \{B \subseteq [m] : |B| = k,\; \{v_i : i \in B\} \text{
  are linearly independent}\}$, where $k = \rank[v_1 \cdots v_m]$.
  The \emph{affine matroid} of hyperplanes with normals
  $w_1, \ldots, w_m \in \R^d$ and biases $b_1, \ldots, b_m \in \R$
  is the linear matroid of the augmented matrix
  $[W \mid b] \in \R^{m \times (d+1)}$, where $W$ has rows $w_i$
  and $b$ has entries $b_i$.
\end{definition}

\begin{definition}[Cyclic interval]\label{def:cyclic-interval}
  A subset $S \subseteq [n] = \{0, 1, \ldots, n{-}1\}$ is a
  \emph{cyclic interval} if $S = \{j, j{+}1, \ldots, j{+}|S|{-}1\}
  \bmod n$ for some $j \in [n]$.  Equivalently, the elements of~$S$
  form a contiguous arc when placed on the circle $\Z/n\Z$.
\end{definition}

\begin{definition}[Grassmann necklace and positroid]
\label{def:grassmann-necklace}
  Let $\cM$ be a matroid of rank~$k$ on~$[n]$.  For $j \in [n]$,
  define the \emph{cyclic order} $\leq_j$ on $[n]$ by
  $j <_j j{+}1 <_j \cdots <_j j{-}1$ (indices mod~$n$).  The
  \emph{Grassmann necklace} of $\cM$ is
  $\cI = (I_0, I_1, \ldots, I_{n-1})$, where $I_j$ is the
  lexicographically smallest basis of $\cM$ in the order $\leq_j$.

  The \emph{Gale order} $\geq_j$ on $k$-subsets is defined by: writing
  $B = \{b_1 <_j \cdots <_j b_k\}$ and
  $I_j = \{i_1 <_j \cdots <_j i_k\}$, we have
  $B \geq_j I_j$ iff $b_\ell \geq_j i_\ell$ for all
  $\ell = 1, \ldots, k$.

  A matroid $\cM$ is a \emph{positroid} if
  $\cB(\cM) = \{B \subseteq [n] : |B| = k,\;
  B \geq_j I_j \text{ for all } j \in [n]\}$.
\end{definition}

\begin{definition}[Totally positive matrix]\label{def:tp}
  A real matrix $M$ is \emph{totally positive (TP)} if every minor
  (determinant of every square submatrix) is strictly positive.  A
  matrix is \emph{totally nonnegative (TN)} if every minor is
  nonnegative.
\end{definition}


\subsection{Network setup}\label{sec:network-setup}

We consider single-hidden-layer ReLU networks
\begin{equation}\label{eq:network}
  f(x) = W_2 \cdot \relu(W_1 x + b_1) + b_2,
\end{equation}
where $W_1 \in \R^{H \times d}$ is the weight matrix, $b_1 \in
\R^H$ is the bias vector, $W_2 \in \R^{c \times H}$ and $b_2 \in
\R^c$ are the output layer parameters, and $\relu(z) = \max(z, 0)$
is applied element-wise.

The $i$-th hidden neuron defines the hyperplane
$h_i = \{x \in \R^d : w_i^\top x + b_i = 0\}$, partitioning input
space into $2^H$ (at most) activation regions
(Figure~\ref{fig:hyperplanes}).  The augmented matrix
$A = [W_1 \mid b_1] \in \R^{H \times (d+1)}$
defines an affine matroid $\cM(A)$ of rank at most $d + 1$ on
ground set $[H]$.

\paragraph{TP weight construction.}
We construct TP weight matrices using Karlin's kernel theory
\citep{karlin1968}: for strictly increasing sequences
$a_1 < \cdots < a_H$ and $b_1 < \cdots < b_d$, the matrix
$(W_1)_{ij} = \exp(a_i \cdot b_j)$ is totally positive.  The
Cauchy kernel $(W_1)_{ij} = 1/(a_i + b_j)$ with strictly
increasing positive sequences provides an alternative TP
construction.  Biases are initialized randomly.

\begin{figure}[t]
  \centering
  \includegraphics[width=0.65\linewidth]{figures/fig-hyperplanes.pdf}
  \caption{Hyperplane arrangement of a TP-constrained ReLU network
    ($H = 8$, $d = 2$) trained on the two-moons dataset.  Each line
    is the decision boundary $w_i^\top x + b_i = 0$ of one hidden
    neuron.  The affine matroid of the augmented matrix
    $[W_1 \mid b_1]$ captures which triples of these hyperplanes
    meet at a common point (non-bases).}
  \label{fig:hyperplanes}
\end{figure}


\subsection{The conjecture}\label{sec:conjecture}

\begin{conjecture}[Activation Positroid Conjecture]
\label{conj:apc}
  If $W_1$ is totally positive, then for any bias vector $b_1$ the
  affine matroid $\cM([W_1 \mid b_1])$ is a positroid.
\end{conjecture}

We tested Conjecture~\ref{conj:apc} across 800+ training trials on
multiple datasets (two-moons, concentric circles, spirals, XOR,
PCA-reduced digits) with hidden dimensions $H \in \{6, 8, 10, 12,
14, 16, 18, 20\}$ and both TP-constrained and unconstrained training.

\begin{center}
\begin{tabular}{lccc}
  \toprule
  \textbf{Mode} & \textbf{Trials} & \textbf{Positroid rate} & \textbf{Non-uniform rate} \\
  \midrule
  TP exponential & 400+ & 100\% & varies \\
  TP Cauchy      & 200+ & 100\% & varies \\
  Unconstrained  & 200+ & 100\% & varies \\
  \bottomrule
\end{tabular}
\end{center}

Every trial produced a positroid activation matroid.  Many are
trivially positroid (uniform matroid $U(d{+}1, H)$, meaning all
hyperplane $(d{+}1)$-tuples are in general position), but
non-uniform matroids---those with genuine dependencies among
hyperplanes---also arise, particularly from the exponential kernel on
the two-moons dataset.

\begin{remark}
  Unconstrained (non-TP) training also produces 100\% positroid rate
  in our experiments, but every unconstrained trial yields a uniform
  matroid---so the positroid property is vacuous.  Whether gradient
  descent itself induces positroid structure when non-uniform matroids
  arise requires a non-TP parameterisation that produces genuine
  dependencies; we investigate this in
  Section~\ref{sec:tp-essential}.
\end{remark}


% ══════════════════════════════════════════════════════════════════════
\section{Disproving the conjecture}\label{sec:counterexamples}
% ══════════════════════════════════════════════════════════════════════

The TP constraint applies to $W_1$ but not to $b_1$.  This gap
allows us to construct counterexamples.

\subsection{Dependency coefficients}\label{sec:dependency}

For a rank-$(d{+}1)$ matroid on $[H]$, each non-basis is a
$(d{+}1)$-element subset $S$ whose corresponding rows of the
augmented matrix $[W_1 \mid b_1]$ are linearly dependent.  This
dependence takes the form
\begin{equation}\label{eq:dependency}
  \sum_{i \in S} c_i \begin{pmatrix} w_i \\ b_i \end{pmatrix} = 0,
\end{equation}
where $c = (c_i)_{i \in S}$ is the dependency coefficient vector,
computable from the left null space of the submatrix indexed by $S$.

When $d = 2$ (rank $k = 3$), each dependency
$c_i [w_i \mid b_i] + c_j [w_j \mid b_j] + c_k [w_k \mid b_k] = 0$
with the TP constraint on $W_1$ forces a specific sign pattern on
$(c_i, c_j, c_k)$: by the TP property, alternating signs $(+, -, +)$
up to global scaling.

Crucially, equation~\eqref{eq:dependency} is \emph{linear} in the
biases $b_i$.  Given fixed TP weights, one can solve for biases that
force any desired subset to become a non-basis.

\subsection{Crossing circuits}\label{sec:crossing}

\begin{definition}[Crossing pair]\label{def:crossing}
  Two subsets $S_1, S_2 \subseteq [n]$ form a \emph{crossing pair}
  if neither $S_1 \setminus S_2$ nor $S_2 \setminus S_1$ can be
  contained in a single arc of the circle~$\Z/n\Z$ between
  consecutive elements of the other set.  Equivalently, $S_1$ and
  $S_2$ ``interleave'' on the circle.
\end{definition}

The key empirical observation is that crossing non-bases are an
obstruction to the positroid property.  If two non-bases form a
crossing pair, the Grassmann necklace reconstruction cannot
simultaneously accommodate both, potentially breaking the positroid
structure.  (We do not prove this in general; the connection is
verified computationally for the matroids in our experiments.)

\subsection{Counterexample construction}\label{sec:construction}

Our \emph{targeted search} employs two strategies:
\begin{enumerate}
  \item Generate a TP weight matrix $W_1$ (exponential kernel).
  \item \textbf{Crossing-pair strategy.}  Identify pairs of
    $(d{+}1)$-subsets that interleave on the circle and solve
    the linear system~\eqref{eq:dependency} for biases $b_1$
    forcing both subsets to become non-bases simultaneously.
  \item \textbf{Single-spread strategy.}  Choose a single
    non-interval $(d{+}1)$-subset and solve for biases making it
    a non-basis.
  \item Verify the resulting affine matroid $\cM([W_1 \mid b_1])$ is
    (a)~TP-weight, (b)~non-uniform, and (c)~non-positroid.
\end{enumerate}
Crossing pairs succeed at ${\sim}89\%$ of non-uniform trials;
single spread non-bases at ${\sim}73\%$.

\begin{example}\label{ex:smallest}
  The smallest counterexample arises at $(d, H) = (2, 5)$: a
  rank-3 matroid on $[5]$ obtained from $U(3, 5)$ by declaring
  $\{0, 2, 4\}$ a non-basis.  The subset $\{0, 2, 4\}$ is \emph{not}
  a cyclic interval (the elements alternate around the circle),
  and the resulting matroid fails the positroid property.
\end{example}

\subsection{Results}\label{sec:counterexample-results}

Across configurations $(d, H) \in \{(2,5), (2,6), (2,8)\}$
with the exponential TP kernel, we constructed
\textbf{12,642 counterexamples}: TP-weight networks with
deliberately chosen biases whose affine matroids are not positroids.

The counterexamples reveal a sharp dichotomy.  Among 1{,}238
non-uniform matroids from 2{,}793 TP network trials (combining
targeted and random bias strategies):

\begin{center}
\begin{tabular}{lcc}
  \toprule
  & \textbf{Positroid} & \textbf{Non-positroid} \\
  \midrule
  All non-bases are cyclic intervals & 433 & \textbf{0} \\
  Has $\geq 1$ non-interval non-basis & 58 & 747 \\
  \bottomrule
\end{tabular}
\end{center}

\noindent The zero in the upper-right cell is not a coincidence---it
is the content of Theorem~\ref{thm:cip} (see also
Figure~\ref{fig:dichotomy}).  The 58 positroids with non-interval
non-bases show that the converse does not hold: having a spread
non-basis is necessary but not sufficient for breaking positroid
structure.

\begin{figure}[t]
  \centering
  \includegraphics[width=0.65\linewidth]{figures/fig-dichotomy.pdf}
  \caption{Counterexample dichotomy.  Among 1{,}238 non-uniform
    matroids from 2{,}793 TP network trials, the zero count (left,
    non-positroid column) is explained by the CIP theorem: cyclic
    interval non-bases \emph{never} break positroid structure.}
  \label{fig:dichotomy}
\end{figure}

\subsection{Revised conjecture}\label{sec:revised-conjecture}

The counterexamples above disprove Conjecture~\ref{conj:apc} as
stated: TP weights alone do not force positroid structure.  However,
the 800+ trained networks (TP and unconstrained) from
Section~\ref{sec:conjecture} all produced positroids, and their
non-basis support is always a contiguous range of neuron
indices---never the gapped or interleaving patterns that the
counterexample search exploits.  This motivates a weaker conjecture:

\begin{conjecture}[Trained Positroid Conjecture]
\label{conj:trained}
  If $W_1$ is totally positive and the network~\eqref{eq:network}
  is trained by gradient descent, then the affine matroid
  $\cM([W_1 \mid b_1])$ is a positroid.
\end{conjecture}

\noindent This shifts the claim from a static algebraic property of
TP matrices to a dynamical one about the interaction between TP
structure and gradient-based optimisation.


% ══════════════════════════════════════════════════════════════════════
\section{The theorems}\label{sec:theorems}
% ══════════════════════════════════════════════════════════════════════

\subsection{Contiguous-Implies-Positroid}\label{sec:cip}

\begin{figure}[t]
  \centering
  \includegraphics[width=0.75\linewidth]{figures/fig-intervals.pdf}
  \caption{Non-basis patterns on the ground set $[6] = \{0,\ldots,5\}$,
    viewed on a circle.  \textbf{Left:} The subset $\{2,3,4\}$ forms a
    contiguous arc (cyclic interval); by the CIP theorem, declaring it a
    non-basis always yields a positroid.  \textbf{Right:} The subset
    $\{0,2,4\}$ is non-contiguous; declaring it a non-basis produces a
    non-positroid (Corollary~\ref{cor:single-removal}).}
  \label{fig:intervals}
\end{figure}

\begin{theorem}[Contiguous-Implies-Positroid (CIP)]
\label{thm:cip}
  Let $\cM$ be a matroid of rank $k$ on $[n]$.  If every non-basis
  of $\cM$ is a cyclic interval (Definition~\ref{def:cyclic-interval}),
  then $\cM$ is a positroid.
\end{theorem}

\begin{proof}
  Let $\cI = (I_0, \ldots, I_{n-1})$ be the Grassmann necklace of
  $\cM$.  By the characterization of positroids
  (Definition~\ref{def:grassmann-necklace}), it suffices to show
  that a $k$-subset $B \subseteq [n]$ is a basis of $\cM$ if and
  only if $B \geq_j I_j$ for all $j \in [n]$.

  \smallskip\noindent
  $(\Rightarrow)$\; If $B$ is a basis, then $B \geq_j I_j$ for
  every $j$, because $I_j$ is by definition the minimum basis in
  the Gale order $\geq_j$.

  \smallskip\noindent
  $(\Leftarrow)$\; Suppose $B$ is \emph{not} a basis.  We exhibit
  $j$ such that $B \not\geq_j I_j$.  By hypothesis, $B$ is a
  cyclic interval, say $B = \{a, a{+}1, \ldots, a{+}k{-}1\} \bmod n$.
  This is the lexicographically \emph{smallest} $k$-subset in the
  cyclic order~$\leq_a$: for any $k$-subset $S$, we have
  $S \geq_a B$.  In particular, $I_a \geq_a B$.  Since $B$ is not a
  basis but $I_a$ is, $I_a \neq B$, whence $I_a >_a B$.  This means
  $B <_a I_a$, so $B \not\geq_a I_a$.
\end{proof}

\begin{corollary}[Single-removal dichotomy]\label{cor:single-removal}
  Let $2 \leq k \leq n - 2$, let $S$ be a $k$-subset of $[n]$, and
  let $\cM = U(k,n) \setminus \{S\}$ be the matroid obtained from
  the uniform matroid by declaring $S$ a
  non-basis.\footnote{More precisely, $\cM$ has bases
  $\binom{[n]}{k} \setminus \{S\}$, which satisfies the exchange
  axiom when $n > k$.}  Then $\cM$ is a positroid if and only if
  $S$ is a cyclic interval (see Figure~\ref{fig:intervals}).
\end{corollary}

\begin{proof}
  The forward direction is Theorem~\ref{thm:cip}.  For the converse,
  suppose $S$ is not a cyclic interval.  Then for every $j$, the
  cyclic interval $\{j, j{+}1, \ldots, j{+}k{-}1\}$ is still a
  basis, so $I_j = \{j, j{+}1, \ldots, j{+}k{-}1\}$.  Since this
  interval is the Gale-minimum $k$-subset in $\leq_j$, every
  $k$-subset Gale-dominates it; in particular $S \geq_j I_j$ for
  all~$j$.  But $S \notin \cB(\cM)$, so $\cM$ has a non-basis that
  satisfies all $n$ Gale conditions and therefore $\cM$ is not a
  positroid.
\end{proof}

\begin{remark}
  The if-and-only-if characterization is specific to single removal.
  For multi-removal matroids (multiple non-bases declared),
  non-interval non-bases can still yield positroids: among $2{,}793$
  TP network trials, $58$ positroids had at least one non-interval
  non-basis.
\end{remark}


\subsection{Contiguous-Support Positroid Theorem}\label{sec:cspt}

The CIP theorem handles the case where individual non-bases are
cyclic intervals.  In practice, trained networks often have
\emph{many} non-bases (at $H = 200$ with $d = 2$, potentially
thousands), but their \emph{support}---the union of all elements
appearing in any non-basis---forms a contiguous arc.

\begin{definition}[Non-basis support]\label{def:support}
  For a matroid $\cM = (E, \cB)$ of rank $k$, the \emph{non-basis
  support} is
  \[
    \supp(\cM) = \bigcup \bigl\{ B \subseteq E : |B| = k,\;
    B \notin \cB \bigr\}.
  \]
  This is the set of elements that participate in at least one
  non-basis.  For the uniform matroid, $\supp(\cM) = \emptyset$.
\end{definition}

\begin{definition}[Support rank deficiency]\label{def:rank-deficiency}
  The \emph{support rank deficiency} of $\cM$ is
  $k - \rank_\cM(\supp(\cM))$, where $\rank_\cM(S)$ is the size of
  the largest independent subset of $S$.
\end{definition}

\begin{theorem}[Contiguous-Support Positroid Theorem (CSPT)]
\label{thm:cspt}
  Let $\cM$ be a matroid of rank $k$ on $[n]$.  Let $S = \supp(\cM)$
  be the non-basis support.  If $S$ is a cyclic interval and
  $\rank_\cM(S) = r < k$, then $\cM$ is a positroid.
\end{theorem}

\begin{proof}
  Since every non-basis $B$ has $B \subseteq S$ by definition of
  support, and $\rank_\cM(S) = r < k = |B|$, every $k$-subset of
  $S$ is dependent, hence a non-basis.  Conversely, any $k$-subset
  \emph{not} entirely contained in $S$ contains an element
  $e \notin S$; if it were a non-basis, then $e$ would be in the
  support, contradicting $e \notin S$.  Therefore:
  \begin{equation}\label{eq:basis-characterization}
    B \in \cB(\cM) \iff B \not\subseteq S.
  \end{equation}

  Let $S = \{a, a{+}1, \ldots, a{+}s{-}1\} \bmod n$ with $s = |S|$,
  and let $\bar S = [n] \setminus S$.  Let $\cI = (I_0, \ldots,
  I_{n-1})$ be the Grassmann necklace.  We verify the Gale order
  characterization.

  \smallskip\noindent
  $(\Rightarrow)$\; Immediate from the definition of $I_j$ as the
  Gale-minimum basis.

  \smallskip\noindent
  $(\Leftarrow)$\; Let $B$ be a non-basis, so $B \subseteq S$ with
  $|B| = k$.  Consider $j = a$.  In the cyclic order~$\leq_a$,
  every element of $S$ precedes every element of $\bar S$:
  $a <_a a{+}1 <_a \cdots <_a a{+}s{-}1 <_a a{+}s <_a \cdots$.

  The lex-min basis $I_a$ is constructed by the greedy algorithm
  starting at $a$.  Since $\rank_\cM(S) = r$, the greedy picks
  exactly $r$ elements from $S$ (the first $r$ that maintain
  independence) before no further element of $S$ can increase rank.
  It then picks $k - r$ elements from~$\bar{S}$.

  Write $I_a = \{i_1 <_a \cdots <_a i_k\}$ and
  $B = \{b_1 <_a \cdots <_a b_k\}$.  The first $r$ elements
  $i_1, \ldots, i_r$ lie in $S$, and $i_{r+1}, \ldots, i_k$ lie
  in $\bar{S}$.  All elements of $B$ lie in $S$.  At position
  $\ell = r + 1$:
  \[
    b_{r+1} \in S, \qquad i_{r+1} \in \bar{S}, \qquad
    b_{r+1} <_a i_{r+1},
  \]
  since $S$ precedes $\bar{S}$ in the order $\leq_a$.  Therefore
  $B \not\geq_a I_a$.
\end{proof}

\begin{corollary}[Tail-collapse positroid]\label{cor:tail-collapse}
  If a trained TP-weight network has non-basis support forming a
  cyclic interval of rank less than $d + 1$, the activation matroid
  is a positroid.
\end{corollary}

\begin{remark}\label{rem:hypotheses-necessary}
  Both hypotheses are necessary.  The cyclic interval condition
  fails for $U(2, 5) \setminus \{\{0, 2\}\}$: the support $\{0, 2\}$
  is not a cyclic interval on $[5]$ (elements $0$ and $2$ are
  separated by~$1$), and the matroid is not a positroid.  The rank
  condition is also essential: if the support were a cyclic interval
  but had full rank~$k$, some $k$-subsets of $S$ would be bases,
  breaking the clean characterization~\eqref{eq:basis-characterization},
  and non-positroid examples can arise.
\end{remark}


% ══════════════════════════════════════════════════════════════════════
\section{TP structure is essential}\label{sec:tp-essential}
% ══════════════════════════════════════════════════════════════════════

The CIP and CSPT theorems explain \emph{why} contiguous support
implies positroid structure, but not \emph{why} trained TP networks
always produce contiguous support.  To isolate the role of total
positivity, we compare eight weight parameterization modes.

\subsection{Parameterization modes}\label{sec:param-modes}

All modes use the same architecture and training procedure (Adam
optimizer, learning rate $0.01$, 200 epochs on the two-moons dataset,
$H \in \{6, 8, 10\}$, $d = 2$).  Weight matrices differ only in
initialization and constraint structure:

\begin{enumerate}
  \item \textbf{TP exponential:} $W_{ij} = \exp(a_i \cdot b_j)$
    with strictly increasing $a, b$.
  \item \textbf{TP Cauchy:} $W_{ij} = 1/(a_i + b_j)$ with
    strictly increasing positive $a, b$.
  \item \textbf{Sinusoidal (non-TP):}
    $W_{ij} = 2 + \sin(a_i \cdot b_j)$.
  \item \textbf{Quadratic distance (non-TP):}
    $W_{ij} = (a_i - b_j)^2 + 1$.
  \item \textbf{Unconstrained:} Xavier/Glorot initialization.
  \item \textbf{Permuted exponential (non-TP):}
    $W_{ij} = \exp(a_i \cdot b_{\pi(j)})$ where $\pi$ reverses the
    column ordering.
  \item \textbf{Negated bidiagonal (non-TP):}
    $W = B \cdot \exp(a b^\top)$ where $B$ is lower bidiagonal with
    alternating-sign subdiagonal entries
    ($B_{i+1,i} = (-1)^{i+1} \cdot 1.5$).
  \item \textbf{Fixed convergent bias-only (non-TP):} Frozen $W_1$,
    train only biases and the output layer.
\end{enumerate}

Modes 1--2 maintain TP structure by parameterizing $a$ and $b$
through cumulative softplus to enforce strict monotonicity during
training.

\subsection{Results}\label{sec:tp-results}

Across 60 moons trials per mode (20 trials $\times$ 3 hidden
dimensions):

\begin{center}
\begin{tabular}{lcc}
  \toprule
  \textbf{Mode} & \textbf{Non-positroid rate} & \textbf{Non-uniform rate} \\
  \midrule
  TP exponential & 0/60 (0\%) & 13/60 (22\%) \\
  TP Cauchy & 0/60 (0\%) & 0/60 (0\%) \\
  Negated bidiagonal & 2/60 (3.3\%) & 15/60 (25\%) \\
  All other non-TP & 0/60 (0\%) & 0/60 (0\%) \\
  \bottomrule
\end{tabular}
\end{center}

\noindent
Note that TP Cauchy and ``all other non-TP'' modes have 0\%
non-uniform rate: every trial produced a uniform matroid, so the 0\%
non-positroid rate is vacuous (all uniform matroids are trivially
positroid).  Only TP exponential and negated bidiagonal produce
non-uniform matroids that genuinely test the positroid question.

The negated bidiagonal mode produced two non-positroid matroids with
gapped (non-contiguous) non-basis supports: $\{1,3,4,5\}$ (at $H=6$,
skipping elements $0$ and $2$) and $\{3,5,6,7,8,9\}$ (at $H=10$,
skipping element $4$).  These gaps violate the hypothesis of the
CSPT.

\subsection{Causal mechanism}\label{sec:causal}

The experimental evidence supports the causal chain:
\begin{equation}\label{eq:causal}
  \underset{\text{(weight structure)}}{\text{TP}}
  \;\longrightarrow\;
  \underset{\text{(training dynamics)}}{\text{contiguous support}}
  \;\longrightarrow\;
  \underset{\text{(Theorem~\ref{thm:cspt})}}{\text{positroid}}
\end{equation}

\begin{itemize}
  \item \textbf{TP $\to$ contiguous support:} The TP property
    constrains the normal vectors to lie in a specific geometric
    configuration.  Empirically, this causes the rank-deficient
    hyperplanes to cluster as a contiguous arc in the cyclic ordering
    inherited from the TP parametrization.

  \item \textbf{Contiguous support $\to$ positroid:} This is the
    content of the CSPT (Theorem~\ref{thm:cspt}).

  \item \textbf{Breaking TP breaks contiguity:} The negated
    bidiagonal mode, which has a controlled non-TP perturbation,
    can produce gapped supports, which the CSPT cannot rescue.
\end{itemize}

\begin{remark}
  The first arrow---why TP weights produce contiguous support under
  gradient descent---remains an open question.  We conjecture it
  relates to the monotone convergence of normal vectors in the TP
  parameterization, where the cyclic ordering of hyperplane normals
  is preserved throughout training.
\end{remark}


% ══════════════════════════════════════════════════════════════════════
\section{Matroid-guided pruning}\label{sec:pruning}
% ══════════════════════════════════════════════════════════════════════

The matroid structure provides a principled criterion for neural
network pruning that does not require retraining, importance scoring,
or gradient computation.

\subsection{Essential and tail neurons}\label{sec:essential-tail}

\begin{definition}[Essential and tail partition]
  Given a network with affine matroid $\cM$ on $[H]$, we define:
  \begin{itemize}
    \item \textbf{Essential neurons:} $[H] \setminus \supp(\cM)$,
      the elements not contained in any non-basis.  Every $k$-subset
      containing an essential neuron is linearly independent.
    \item \textbf{Tail neurons:} $\supp(\cM)$, the elements that
      participate in at least one non-basis.
  \end{itemize}
\end{definition}

At input dimension $d = 2$, the affine matroid has rank $k = 3$, so
at most 3 neurons carry independent geometric information.  At
$H = 200$, over 170 neurons are tail---they contribute no
additional combinatorial structure to the hyperplane arrangement
beyond what the essential neurons provide
(Figure~\ref{fig:pruning}a).

\subsection{Pruning strategies}\label{sec:strategies}

We compare five pruning strategies:

\begin{enumerate}
  \item \textbf{Matroid-guided:} Remove tail neurons
    entirely---delete the corresponding rows of $W_1$, entries of
    $b_1$, and columns of $W_2$.  The network shrinks from $H$ to
    $H - |\text{pruned}|$ neurons.

  \item \textbf{Random:} Remove the same number of neurons as
    matroid-guided but selected uniformly at random.  Averaged over
    20 draws.

  \item \textbf{Magnitude:} Remove neurons with smallest $\ell^2$
    norm $\|w_i\|_2$ of their weight vectors.

  \item \textbf{Activation:} Remove neurons with lowest mean ReLU
    activation $\mathbb{E}_x[\relu(w_i^\top x + b_i)]$, analogous
    to the APoZ criterion.

  \item \textbf{Sensitivity:} Remove neurons with smallest
    importance score $|W_{2,i}| \cdot \|w_i\|_2$, combining output
    weight magnitude with input weight norm.
\end{enumerate}

\subsection{Results}\label{sec:pruning-results}

We train TP-exponential networks on the two-moons dataset at
$H \in \{50, 100, 200\}$ with $d = 2$, 200 training samples,
and evaluate pruning at fractions $\{0\%, 25\%, 50\%, 75\%, 100\%\}$
of tail neurons (Figure~\ref{fig:pruning}).

\begin{figure}[t]
  \centering
  \includegraphics[width=\linewidth]{figures/fig-pruning.pdf}
  \caption{\textbf{(a)}~Essential (thick) vs.\ tail (thin) hyperplanes
    in a trained TP network ($H = 20$, $d = 2$).  Only 4 essential
    neurons carry independent geometric information; the remaining 16
    are redundant in the affine matroid.
    \textbf{(b)}~Pruning curves at $H = 200$.  Matroid-guided pruning
    maintains zero accuracy loss through 75\% tail removal, outperforming
    random pruning by 14.4 percentage points.}
  \label{fig:pruning}
\end{figure}

\paragraph{75\% tail removal: zero accuracy loss.}
Table~\ref{tab:pruning-75} shows the accuracy delta (pruned minus
original) at 75\% tail removal across all scales and strategies.

\begin{table}[t]
  \centering
  \caption{Accuracy delta at 75\% tail removal.  Matroid-guided
    achieves zero loss at every scale; random pruning degrades
    by 10--15 percentage points.  On this dataset, magnitude,
    activation, and sensitivity heuristics select identical neuron
    sets.}
  \label{tab:pruning-75}
  \begin{tabular}{lccc}
    \toprule
    \textbf{Strategy} & $H=50$ & $H=100$ & $H=200$ \\
    \midrule
    Matroid-guided & $0.0\%$ & $0.0\%$ & $0.0\%$ \\
    Magnitude & $-0.7\%$ & $-0.5\%$ & $-0.3\%$ \\
    Activation & $-0.7\%$ & $-0.5\%$ & $-0.3\%$ \\
    Sensitivity & $-0.7\%$ & $-0.5\%$ & $-0.3\%$ \\
    Random & $-10.2\%$ & $-15.4\%$ & $-14.4\%$ \\
    \bottomrule
  \end{tabular}
\end{table}

\paragraph{Pruning curves.}
At $H = 200$, the matroid-guided strategy maintains 0\% accuracy
delta through 75\% tail removal, dropping to $-0.7\%$ only at 100\%
(removing all tail neurons).  Random pruning degrades monotonically,
reaching $-20.9\%$ at 100\%.  Heuristic baselines (magnitude,
activation, sensitivity) plateau at $-0.3\%$.  The three heuristics
select identical neuron sets because the exponential kernel assigns
monotonically increasing weight norms to higher-index neurons, so
magnitude, mean activation, and gradient sensitivity all produce the
same pruning order---this is a property of the TP parameterization,
not evidence that the methods are equivalent in general.
The advantage of matroid-guided over random is approximately 14.4
percentage points at 75\% pruning.

\paragraph{Scaling behavior.}
The advantage of matroid-guided pruning is consistent across scales:
10--15pp over random at every $H$ tested.  The absolute number of
pruned neurons grows (from $\sim$29 at $H = 50$ to $\sim$138 at
$H = 200$), but the accuracy impact remains zero through 75\%.

\subsection{Why it works}\label{sec:why-pruning}

The matroid partition identifies functional redundancy at the level
of the hyperplane arrangement.  Tail neurons define hyperplanes that
are \emph{linearly dependent} on the essential neurons' hyperplanes
(in the affine sense).  Removing a tail neuron does not eliminate any
activation region boundary that isn't already determined by the
essential neurons.

More precisely, let $A = [W_1 \mid b_1]$ be the augmented matrix.
When $H \geq 2(d+1)$ (as in all our experiments), the essential set
$[H] \setminus \supp(\cM)$ contains at least $d + 1$ elements whose
rows span the full row space of $A$.  Tail neuron rows lie in this
span and are therefore redundant.

\begin{remark}[100\% tail removal]
  Removing \emph{all} tail neurons is initialization-dependent.  At
  $H = 200$, 100\% tail removal causes $-0.7\%$ accuracy loss---small
  but nonzero.  This is because the network's output layer weights
  $W_2$ assign nonzero weight to tail neurons, and removing them
  changes the network function even though the arrangement
  combinatorics are preserved.  At 75\%, the remaining tail neurons
  absorb enough of this output weight to maintain accuracy.
\end{remark}


% ══════════════════════════════════════════════════════════════════════
\section{Discussion}\label{sec:discussion}
% ══════════════════════════════════════════════════════════════════════

\subsection{Limitations}\label{sec:limitations}

\paragraph{Input dimension.}
All experiments use $d = 2$ (rank-3 affine matroids).  The theorems
(CIP, CSPT) are dimension-agnostic, but the empirical verification
is limited to rank 3.  At $d \geq 3$, the exponential TP kernel
has only $H + d$ parameters---too few degrees of freedom to span
$\mathbb{R}^d$, confining normals to a one-dimensional curve---and
alternative TP constructions (Loewner--Whitney) produce uniform
matroids.  Extending the empirical investigation to
higher rank is an important direction.

\paragraph{Single hidden layer.}
We study only the first hidden layer's hyperplane arrangement.
Multi-layer networks have more complex activation region structures
that are not captured by a single affine matroid.

\paragraph{Dataset simplicity.}
The two-moons dataset is a simple binary classification task.  The
pruning results on this dataset demonstrate the principle but do not
directly predict performance on high-dimensional tasks.  The
baselines (magnitude, activation, sensitivity) may have stronger
relative performance on more complex datasets.

\subsection{Connections to positive geometry}\label{sec:positive-geometry}

The appearance of positroids in neural networks is tantalizing from
the perspective of positive geometry \citep{arkani-hamed2017}.  The
totally nonnegative Grassmannian $\Gr_{\geq 0}(k,n)$ and its
positroid stratification play a central role in the amplituhedron
program for scattering amplitudes.  Whether the connection between
neural network training and positive geometry extends beyond the
observations in this paper remains speculative.

\subsection{Open questions}\label{sec:open}

\begin{enumerate}
  \item \textbf{Why does TP $+$ gradient descent produce contiguous
    support?}  The first arrow in the causal chain~\eqref{eq:causal}
    is empirical.  A proof would likely need to analyze the gradient
    flow on the TP manifold and show that normal vector convergence
    preserves contiguity.

  \item \textbf{Higher rank.}  Do the CIP and CSPT phenomena persist
    at rank $\geq 4$?  Computational challenges (enumerating all
    $\binom{H}{k}$ subsets) limit direct verification, but proxy
    metrics (contiguous window ranks, random sampling) could extend
    the investigation.

  \item \textbf{Multi-removal characterization.}  Among the 58
    positroids with non-interval non-bases in our experiments, what
    additional structure characterizes when a set of spread non-bases
    is compatible with the positroid property?

  \item \textbf{Multi-layer positroid signature.}  Can the positroid
    analysis be extended across layers, perhaps via a sequence of
    matroids or a matroid on the full network?

  \item \textbf{Unconstrained training.}  Unconstrained networks
    also achieve 100\% positroid rate.  Is gradient descent on
    \emph{any} architecture an implicit TP regularizer, or does this
    depend on specific features of SGD (momentum, learning rate,
    batch size)?
\end{enumerate}


% ══════════════════════════════════════════════════════════════════════
\section{Conclusion}\label{sec:conclusion}
% ══════════════════════════════════════════════════════════════════════

We have established a new connection between positroid combinatorics
and the geometry of trained ReLU networks.  The
Contiguous-Implies-Positroid theorem and its generalization (CSPT)
explain why trained networks with TP weights always produce positroid
activation matroids: training dynamics produce contiguous non-basis
support, and contiguous support forces positroid structure.
Counterexamples with adversarial biases (12{,}642 instances) show
that the conjecture fails without training, precisely when non-bases
exhibit spread (non-contiguous) patterns.

The matroid structure is not merely a theoretical curiosity.  The
essential/tail partition provides a principled pruning criterion that
removes 75\% of tail neurons with zero accuracy loss, outperforming
four standard baselines.  This suggests that matroid-theoretic analysis of
neural network geometry can yield practical algorithmic benefits.

\paragraph{Code availability.}
Source code and experiments are available at
\url{https://github.com/HarrisonTotty/positroid-structure-relu-networks}.


% ══════════════════════════════════════════════════════════════════════
% References
% ══════════════════════════════════════════════════════════════════════

\bibliographystyle{plainnat}

\clearpage
\begin{thebibliography}{20}

\bibitem[Arkani-Hamed et~al.(2018)]{arkani-hamed2017}
N.~Arkani-Hamed, Y.~Bai, S.~He, and G.~Yan.
\newblock Scattering forms and the positive geometry of kinematics,
  color and the worldsheet.
\newblock \emph{Journal of High Energy Physics}, 2018(5):096, 2018.

\bibitem[Gunasekar et~al.(2018)]{gunasekar2018}
S.~Gunasekar, J.~Lee, D.~Soudry, and N.~Srebro.
\newblock Implicit bias of gradient descent on linear convolutional
  networks.
\newblock In \emph{NeurIPS}, 2018.

\bibitem[Han et~al.(2015)]{han2015}
S.~Han, J.~Pool, J.~Tran, and W.~Dally.
\newblock Learning both weights and connections for efficient neural
  networks.
\newblock In \emph{NeurIPS}, 2015.

\bibitem[Hanin and Rolnick(2019)]{hanin2019}
B.~Hanin and D.~Rolnick.
\newblock Complexity of linear regions in deep neural networks.
\newblock In \emph{ICML}, 2019.

\bibitem[Karlin(1968)]{karlin1968}
S.~Karlin.
\newblock \emph{Total Positivity}.
\newblock Stanford University Press, 1968.

\bibitem[Knutson et~al.(2013)]{knutson2013}
A.~Knutson, T.~Lam, and D.~Speyer.
\newblock Positroid varieties: juggling and geometry.
\newblock \emph{Compositio Mathematica}, 149(10):1710--1752, 2013.

\bibitem[{\L}ukowski et~al.(2023)]{lukowski2023}
T.~{\L}ukowski, M.~Parisi, and L.~Williams.
\newblock The positive tropical Grassmannian, the hypersimplex, and
  the $m=2$ amplituhedron.
\newblock \emph{International Mathematics Research Notices},
  2023(19):16778--16836, 2023.

\bibitem[Oh(2011)]{oh2011}
S.~Oh.
\newblock Positroids and Schubert matroids.
\newblock \emph{Journal of Combinatorial Theory, Series A},
  118(8):2426--2435, 2011.

\bibitem[Postnikov(2006)]{postnikov2006}
A.~Postnikov.
\newblock Total positivity, Grassmannians, and networks.
\newblock \emph{arXiv:math/0609764}, 2006.

\bibitem[Soudry et~al.(2018)]{soudry2018}
D.~Soudry, E.~Hoffer, M.~S. Nacson, S.~Gunasekar, and N.~Srebro.
\newblock The implicit bias of gradient descent on separable data.
\newblock \emph{Journal of Machine Learning Research}, 19(70):1--57,
  2018.

\bibitem[Zhang et~al.(2018)]{zhang2018}
L.~Zhang, G.~Naitzat, and L.-H. Lim.
\newblock Tropical geometry of deep neural networks.
\newblock In \emph{ICML}, 2018.

\end{thebibliography}

\end{document}
